\section{导读：为什么 2025 年又开始“雕 Attention”}

本节先建立整期的技术背景。杨松琳讨论的是大模型基础架构里最核心的模块之一：Attention。过去几年，Transformer 的 FFN 部分已经被 MoE（Mixture of Experts，混合专家）大幅改造，DeepSeek 等模型让 MoE 成为共识；下一块可能被“雕”的，就是 Attention。原因很直接：长上下文、长 CoT、Agentic AI 和长文本 decoding，把 Full Attention 的计算与 KV cache 压力推到了前台。

本期要回答四个问题。第一，Linear Attention 的 “linear” 到底线性在哪里。第二，Kimi Linear 的 KDA 模块为什么重要。第三，Kimi 的混合线性注意力、DeepSeek 的稀疏注意力、MiniMax M2 回到 Full Attention，分别意味着什么取舍。第四，为什么架构研究不能只讲数学优雅，还必须和硬件亲和、并行算法、矩阵乘和 kernel 优化协同。

\begin{importantbox}{本期核心命题}
当数据增长变慢、长上下文需求上升、推理成本变重时，架构创新会重新变得关键。Attention 的新路线不是为了替换 Transformer 这个名字，而是为了在相同 FLOPs 下获得更低 loss、更低 KV cache 压力和更好的长序列 decoding 效率。
\end{importantbox}

\begin{knowledgebox}{视觉策略说明}
本视频是固定访谈画面，没有 slides 或论文投屏。正文不重复插入人物帧；所有图均为自制概念图，用来解释注意力复杂度、KDA、三条 Attention 路线、Scaling Ladder、算法考古和硬件亲和。
\end{knowledgebox}

\subsection{本章小结}

EP119 是架构综述而不是普通访谈。它把 2025 年模型公司在 Attention 上的不同押注放在同一张地图中：Kimi 走 Linear/Hybrid，DeepSeek 走 Sparse，MiniMax M2 暂时回到 Full Attention。

